\section{A quartic in 498 variables}\label{sec:v12}

We now reduce the 523-variable construction of
Theorem~\ref{thm:v9-quartics} to 498 variables.  The input remains one
rational variable $X$; every other variable remains existentially
quantified over $\Q$.  The changes below use the same normalized curve,
the same integer labels, the same three height detectors, and the same
published and manuscript inputs as Section~\ref{sec:v9}.  In particular,
they do not replace the restricted converse, the exponent-nineteen
height theorem, or the simultaneous-prime argument by new assumptions.

\begin{theorem}\label{thm:v12-quartic}
There is a specified integer polynomial $P_{498}$ of total degree four
in 498 variables such that
\[
 X\in\Z\quad\Longleftrightarrow\quad
 \exists\mathbf v\in\Q^{497}:P_{498}(X,\mathbf v)=0.
\]
The polynomial has 24,669 nonzero monomials, content one, and largest
absolute coefficient $16\cdot3^{2^{26}}$.  Its symbolic coefficient
table uses the fixed integer $C=3^{2^{25}}$, not an additional variable.
There is also a specified integer quartic $P_{514}$ with
\[
 X\in\Z\quad\Longleftrightarrow\quad
 \exists\mathbf w\in\Q^{513}:P_{514}(X,\mathbf w)=0.
\]
It has 24,716 nonzero monomials, content one, and largest absolute
coefficient $3^{2048}+1005072314676622<2^{3247}$.
\end{theorem}

The proof distinguishes two sorts of changes.  Some are rational
changes of private witnesses with a rational inverse.  Others alter
the arithmetic certificate while preserving both soundness and the
existence of witnesses for every integer input.  For the latter we
prove those two assertions separately; we do not claim equality of
all intermediate witness assignments.

\subsection{Conventions and retained inputs}\label{subsec:v12-inputs}

Throughout put
\[
 K=\Q(\rho),\quad \rho^2=2,\quad
 d_*=3-\rho,\quad D=d_*^2=11-6\rho,\quad
 E:y^2=x^3-Dx.
\]
The point $P=P_*$ is the same normalized generator as in
Section~\ref{subsec:v5-point-source} and Section~\ref{sec:v9}:
\[
 P=(c_*,a_*b_*),\qquad
 a_*=-1+\tfrac74\rho,\quad
 b_*=\tfrac{39}{8}-\tfrac92\rho,\quad
 c_*=a_*^2=\tfrac{57}{8}-\tfrac72\rho.
\]
Thus $b_*^2=c_*^2-D$.  In particular $a_*,b_*,c_*,D$ are nonzero.
We use the inherited equality $2E(K)=\Z P$, the nonzero point domains
of Lemma~\ref{lem:v5-point-domains}, and the inherited estimate
$\operatorname h(x([n]P))\ge n^2-6$.  Here $\operatorname h$ is the
absolute logarithmic Weil height.  At every unramified place above a
rational prime $p$ we normalize $v(p)=1$.  The coefficient basis is
always $(1,\rho)$ and conjugation sends $\rho$ to $-\rho$.

Retain the four odd-coset root charts of
Proposition~\ref{v9:prop:odd-chart}, with labels
\[
 Q_1=[n]P,\quad Q_2=S=[a]P,\quad
 Q_3=[t]P,\quad Q_4=[b]P,
 \qquad n,a,t,b\text{ odd}.
\]
The letter $b$ in this display is an integer point label.  We use
$\mathfrak b$ below for a rational input to a pole predicate.
Put $Q_0=P$, $U_0=H=[h]P$, $U_i=[u_i]P$, and
$V_i=U_i-Q_i$.  The scalar point is
\[
 R=S-P=[s]P,\qquad s=a-1,
\]
and the input label is $\beta=2X^2+3$.

We retain the pole obstruction and pole-completeness statements
identified in Section~\ref{subsec:v9-inputs}, specifically
Lemmas~\ref{hyp:o8-pole}, \ref{lem:original-seven-simultaneous-realization},
and \ref{lem:v5-one-place}, including their existing source dependencies.
We also retain the three scales of Corollary~\ref{v9:cor:scales} and
its product-height estimate.  Write $\mathcal H=H_{19}$, so that
\begin{equation}\label{eq:v12-height-input}
 1\le\mathcal H<2^{2093966},\qquad
 B^2-6<7\mathcal H|N|^{19/k}
\end{equation}
whenever the three contact radicals divide a nonzero integer $N$
to order $k$, and $B=\max(|n|,|a|,|t|,|b|)$.
This is exactly \eqref{v9:eq:height-transfer}; none of its scales is
replaced by a field norm in the present construction.

The local formal-group inputs are the integral formal group in
$z=-x/y$, its multiplication series, and its logarithm.  The integral
construction is the one in \cite[IV\S1, Example~IV.2.2.3 and
Proposition~IV.2.3(a)]{silverman}.  At the unramified places used below,
$p>76176$, and the logarithm is an isometry on depth at least one by
\cite[Theorem IV.6.4, p.~132]{silverman}; its hypothesis is
$1>1/(p-1)$.  The coefficient and truncation estimates are those of
Lemma~\ref{lem:quartic-reciprocal-log} and its unramified extension in
Lemma~\ref{lem:v5-unramified-log}.  We spell out their new applications
below.  No minimal-model assertion at the prime above two is needed.

\subsection{Mixed logarithmic precision}\label{subsec:v12-mixed}

Let $f=1/x$ and $z=-x/y$, as before.  Define
\[
 a_0=1,\quad a_1=-\tfrac25,\quad
 a_j=\frac{4j-5}{4j+1}a_{j-1}\quad(j\ge2),\qquad
 G_J(T)=\sum_{j=0}^J a_jT^j.
\]
Thus the source truncations are
$\ell_{12}=zG_2(Df^2)$ and $\ell_{40}=zG_9(Df^2)$.
Keep every primary logarithm equal to the uncorrected parameter $z$.
Only the alternate numerator for $V_1$ is shortened:
\begin{equation}\label{eq:v12-ratios}
 L_i=\frac{z(U_i)}{z(H)},\qquad
 L'_1=\overline{\frac{\ell_{12}(V_1)}{\ell_{40}(V_0)}},\qquad
 L'_i=\overline{\frac{\ell_{40}(V_i)}{\ell_{40}(V_0)}}\quad(2\le i\le4).
\end{equation}
The only denominator logarithm remains $\ell_{40}(V_0)$, with its
original nonzero domain.  In particular this change introduces no
new logarithm denominator.

The sixteen field tests and their orders are now
\begin{equation}\label{eq:v12-tests}
\begin{array}{c|l}
1&f(U_1),\ldots,f(U_4),\ \overline{f(V_0)},\ldots,\overline{f(V_4)}\\[2pt]
12&L_1-L'_1,\ L_1-\beta\\[2pt]
40&f(H),\ L_2-L'_2,\ L_3-L'_3,\ L_4-L'_4,\ F(L_2,L_3,L_4),
\end{array}
\end{equation}
where
\[
 F(A,T,B)=T^2-A(A-2)TB+(A-2)^2B^2-1.
\]
This changes the two order-forty input tests in the 523-variable
source to order twelve.  The three other label comparisons and the
Lucas test remain of order forty.

\begin{lemma}\label{lem:v12-mixed-local}
At every consumed unramified place above $p>76176$, the tests
\eqref{eq:v12-tests} imply
\[
 L_1\equiv n\pmod {p^{12}},\qquad
 L_i\equiv n_i\pmod {p^{40}}\quad(2\le i\le4),
 \qquad(n_2,n_3,n_4)=(a,t,b).
\]
The primary points $U_i$ have depth at least twenty after integrality
of the ratios has been established.
\end{lemma}
\begin{proof}
The order-one tests put the indicated points in the formal group.
The test on $f(H)$ gives $v(z(H))\ge20$.  For a formal point $Q$ of
depth $r\ge1$, the unchanged logarithm expansion gives
\[
 \log_E(Q)/z(Q)\in1+p^{4r}\mathcal O.
\]
Its truncation after $G_2$ has relative error of order at least twelve,
and its truncation after $G_9$ has relative error of order at least
forty.  To check these particular applications, the first omitted
indices are respectively $j=3$ and $j=10$ in the source recurrence.
For $4j+1<p$ the coefficients are integral.  For the remaining indices
the source factorial bound
$v_p(a_j)\ge-(4j+1)/(p-1)$ is dominated by $4jr$ and gives the same
stated lower bounds.  Both truncating factors are units.

The full ratios are the rational numbers
\[
 r_i=\frac{u_i}{h},\qquad r'_i=\frac{u_i-n_i}{h-1}.
\]
At least one of $h,h-1$ is a $p$-unit, so at least one of these two
ratios is integral.  Before any improved depth is inferred for $U_i$,
the unit factors just described show
$v(L_i)=v(r_i)$ and $v(L'_i)=v(r'_i)$.  Each cross-test has positive
order, so integrality transfers to both approximate ratios and both
full ratios.  This first step is essential.

Now $L_i$ integral gives $v(z(U_i))\ge v(z(H))\ge20$.
The primary relative errors have order at least eighty.  The alternate
relative errors have order twelve for $i=1$ and forty for $i\ge2$:
in the first case the denominator still has order forty, so it does
not reduce the numerator's order twelve.  Since the ratios are
integral, these relative errors give absolute congruences.  Together
with the cross-tests they yield $L_i\equiv r_i\equiv r'_i$ to the
asserted order.  Multiply by $h$ and $h-1$ and subtract.  The result
is $L_i\equiv n_i$, without division by either integer.  At an inert
prime conjugation preserves the single local valuation, so this
argument has the same conclusion there.
\end{proof}

For clarity, the detector's product is unchanged as a mathematical
operation.  For a test $D_j$ of order $m_j$, a scale $\mathcal A$, and
$t=\mathcal A\lambda$, $\delta=t-1$, its factors are
\[
 1+\mathcal R\left(\frac{D_j}{\delta^{m_j+1}\mathcal R^{m_j}}\right)^2,
 \qquad
 \mathcal R=\mathcal A^{-5}F_5(t)/\delta^6,
 \qquad F_5(t)=\prod_{\theta\in\Theta}(t-\theta),
\]
with the same $\Theta$ as \eqref{eq:v5-detector-data}.  On writing
$e=(\mathcal A\delta)^5/F_5(t)$, the multiplier of $D_j^2$ is
$q_{m_j}=q_1e^{2m_j-2}$, where $q_1=e/\delta^3$.
For each $m\in\{1,12,40\}$, form the coefficients of
\[
 \prod_{j:m_j=m}(1+TD_j^2)
\]
and evaluate at $T=q_m$.  Their degrees are respectively $9,2,5$.
Their product is exactly the same sixteen-factor detector expression.
The smaller $G_2$ numerator and these three grouped products explain
the changed quadratic circuit.  Multiplying in $K$ means the two
rational bilinear identities in the fixed basis.  Every reassociation
or grouping here follows from distributivity, with no division or
new witness domain.  The retained coefficient tables and exact
quadratic elimination identities specify the resulting circuit.

\subsection{Five rational half-trace guards}\label{subsec:v12-guards}

For each $0\le i\le4$ write the subtraction denominator as
\[
 d_i=x(U_i)-x(Q_i)=A_i+B_i\rho.
\]
The old field inverse used two private rational coordinates.  Replace
those two equations by the single equation
\begin{equation}\label{eq:v12-half-trace}
 A_i\gamma_i=1,\qquad\gamma_i\in\Q.
\end{equation}
No other equation used the removed inverse coordinates.  This saves
one coordinate in each of five blocks.  It strengthens the domain;
there is no claim that every old trace-zero assignment is retained.

Soundness has an explicit rational map.  From \eqref{eq:v12-half-trace}
we have $A_i\ne0$, hence $d_i\ne0$ and
$A_i^2-2B_i^2\ne0$.  Restore the old field inverse as
\[
 d_i^{-1}=\frac{A_i-B_i\rho}{A_i^2-2B_i^2}.
\]
All old subtraction outputs and all their consumers are unchanged.

For completeness fix the four root points and the scalar preparation
before choosing contact primes, as in Section~\ref{subsec:v9-completeness}.
Increase the allowed finite prime cutoff to exclude every prime at
which any of $P,Q_1,\ldots,Q_4$ has nonintegral abscissa.  This is a
finite restriction on already fixed points.  Choose any one of the
split contact primes supplied by Lemma~\ref{lem:v5-two-reductions},
and denote its primary and opposite places by $v,\bar v$.
The unchanged CRT construction makes every $U_i$ formal at $v$ and
makes it reduce to $Q_i$ at $\bar v$.  All five $Q_i$ have finite
integral reduction.  Consequently
\[
 v(d_i)<0,\qquad \bar v(d_i)\ge0.
\]
Conjugation exchanges the two places, so
$v(\overline{d_i})=\bar v(d_i)\ge0$.  The two terms of
$d_i+\overline{d_i}=2A_i$ have unequal valuations; the sum is nonzero.
Thus all five inverses $\gamma_i=1/A_i$ exist for the same completeness
witnesses.  This includes $i=0$.  It imposes no unit condition on the
root labels and no condition on the permitted zero pole witnesses.

\subsection{Normalizing the three detector poles}\label{subsec:v12-poles}

Here $\mathfrak b$ is one detector's rational pole input.  In the
compiler's rational coordinates the seven relevant equations are
\begin{align*}
 p^2&=p_2,&17\mathfrak b p_2+3e^2-34e\mathfrak b-289\mathfrak b^2&=0,\\
 u_2&=\mathfrak b^2,&17e\mathfrak b-eu_2-17g&=0,\\
 s_0^2&=s_2,&q_0^2&=q_2,\\
 &&17\mathfrak b q_2+867g^2-34gs_2-s_2^2&=0.
\end{align*}
In particular the stored $g$ is the source's quantity $Bh/17$,
not $Bh$.  The following substitution respects that factor:
\[
 e=\mathfrak b E_0,\quad p=\mathfrak b P_0,\quad
 s_0=\mathfrak b T_0,\quad q_0=\mathfrak b^2Q_0,
\]
\[
 p_2=\mathfrak b^2P_2,\quad s_2=\mathfrak b^2T_2,\quad
 q_2=\mathfrak b^4Q_2,\quad
 u_2=\mathfrak b^2,\quad g=\mathfrak b^2G_0.
\]
The replacement consists of six quadratic equations:
\begin{align}\label{eq:v12-pole-normalized}
 P_0^2&=P_2,&17\mathfrak b P_2+3E_0^2-34E_0-289&=0,\notag\\
 17E_0-E_0\mathfrak b-17G_0&=0,&T_0^2&=T_2,\\
 Q_0^2&=Q_2,&17\mathfrak b Q_2+867G_0^2-34G_0T_2-T_2^2&=0.\notag
\end{align}

The division by $\mathfrak b$ in the reverse coordinate change is
legitimate in both systems.  Indeed the unchanged detector equations,
with its actual product $W=w_0+w_1\rho$, include
\[
 w_0\iota=1,\qquad
 2w_0\nu=w_0^2-2w_1^2,
\]
and an equation whose specialization at $\mathfrak b=0$ is $\nu^2=0$.
Thus $\mathfrak b=0$ would give $\nu=0$ and $w_0^2=2w_1^2$.
Rationality then forces $w_0=w_1=0$, contradicting $w_0\iota=1$.

Under the displayed substitution the old $u_2$ equation is an identity.
The old $p_2$, first-conic, $g$, and $s_2$ equations are
$\mathfrak b^2$ times their replacements; the old $q_2$ and second-conic
equations are $\mathfrak b^4$ times their replacements.  These identities
prove mutually inverse rational maps fixing every outside coordinate.
They remain valid when any of $e,p,s_0,q_0,g$ is zero.  Applying the
change to three detectors saves three coordinates.  The scalar pole
with its preparation variable is not altered by this normalization.

\subsection{One coefficient in the physical anchor}\label{subsec:v12-anchor}

Put $C=3^{2^{25}}$ and write the actual scalar value as
$f(R)=f_0+f_1\rho$.  The previous fixed parameter was $S_0=C^2$.
The old three-witness anchor encoded
$f_0/C,f_1/C\in\Z_{(2)}\cap\Z_{(3)}$.  On the present integer orbit it
is enough to impose that condition on $f_1/C$.

\begin{lemma}\label{lem:v12-anchor-coefficient}
For every nonzero even integer $s$, write
$f([s]P)=f_0+f_1\rho$.  Both coefficients are $2$-integral, and
\[
 v_3(f_1)=2+2v_3(s),\qquad
 v_3(f_0)\ge3+2v_3(s).
\]
In particular $f_1\ne0$, and
\[
 f_1/C\in\Z_{(2)}\cap\Z_{(3)}
 \quad\Longleftrightarrow\quad
 f_0/C,f_1/C\in\Z_{(2)}\cap\Z_{(3)}.
\]
\end{lemma}
\begin{proof}
At $3$ the field is unramified quadratic with residue field
$\F_9=\F_3[\rho]/(\rho^2-2)$.  The exact starting calculation in
Section~\ref{subsec:o8-anchor-subgroup} is
\[
 z(2P)=
 \frac{3396356159088-4135204317756\rho}{3287811625151},
 \qquad z(2P)/3\equiv2+\rho\pmod3.
\]
The second congruence follows by reducing the displayed rational
coefficients after division by three.  Its square is $\rho$ in
$\F_9$.  Write $s=2\cdot3^k u$ with $3\nmid u$, allowing negative $u$.
The integral formal group has inverse $-z$, and hence odd
multiplication series.  Thus
\[
 [u]z=uz+O(z^3),\qquad [3]z=3z+O(z^3)
\]
with integral coefficients.  At depth $r\ge1$, division of the second
identity by $3^{r+1}$ leaves an error of valuation at least $2r-1>0$.
Multiplication by three therefore preserves the leading residue.
Multiplication by $u$ multiplies it by the nonzero residue of $u$.
Consequently
\[
 z([s]P)/3^{k+1}\equiv u(2+\rho)\pmod3.
\]
The exact curve relation is $f=z^2(1-Df^2)$.  On the formal branch its
second factor is one modulo three, and $u^2=1$ in $\F_3$.  Therefore
\[
 f([s]P)/3^{2k+2}\equiv\rho\pmod3.
\]
Independence of $1,\rho$ over $\F_3$ proves the two asserted
coefficient valuations, including the possible zero value of $f_0$.

At the place above two normalize $v(2)=1$.  The literal point has
$v(x(P))=-3$, $v(y(P))=-9/2$, and $v(z(P))=3/2$.  It lies in the
integral formal subgroup, which contains all its integral multiples.
One may construct that subgroup directly from
$w=z^3-Dzw^2$, where $w=-1/y$, as in
Section~\ref{subsec:o8-anchor-subgroup}; this uses no assertion that
the displayed equation is minimal at two.  Thus $f([s]P)$ is integral.
The equality
\[
 v(a+b\rho)=\min\{v_2(a),v_2(b)+\tfrac12\}\quad(a,b\in\Q)
\]
shows that both rational coefficients are $2$-integral.
Since $C$ is a $2$-unit, the equivalence now follows from the
coefficient valuations at three.
\end{proof}

The three-witness criterion in the proof of
Lemma~\ref{lem:v5-three-witness-anchor} applies to every rational $t$:
\begin{equation}\label{eq:v12-quaternion-criterion}
 t\in\Z_{(2)}\cap\Z_{(3)}
 \quad\Longleftrightarrow\quad
 \exists u,v,w\in\Q:\ u^2+6v^2-2w^2-3t^2+2=0.
\end{equation}
To recall its exact input, $(2,3)_\Q$ ramifies precisely at two and
three.  The embedding criterion
\cite[Proposition 14.6.7, pp.~223--224]{voight-quaternion} tests the nonsquareness of
$1-2/(3t^2)$ at those places when $t\ne0$.  At two, positive valuation
of $t$ gives odd negative valuation of the expression, valuation zero
gives residue three modulo eight, and negative valuation gives a
square in $1+8\Z_2$.  At three, nonnegative valuation gives odd negative
valuation of the expression, whereas negative valuation gives a square
in $1+3\Z_3$.  The expression is nonzero because $t^2=2/3$ has no
rational solution.  Multiplication by $3t^2$ gives
\eqref{eq:v12-quaternion-criterion}; for $t=0$ the witnesses are
$(0,0,1)$.  These are exactly the field, local conditions, and
exceptional case already present in that source proof.

Apply this criterion to $t=f_1/C$, multiply its witnesses by $C$,
and then replace the old $w$ by $w+C$.  The resulting anchor is
\begin{equation}\label{eq:v12-anchor}
 u^2+6v^2-2w^2-4Cw-3f_1^2=0.
\end{equation}
It uses the same three private rational witnesses and removes the
old coordinate for $f_0^2+f_0f_1-f_1^2$.  Since $a$ is odd and $R$ is
nonidentity, $s=a-1$ is a nonzero even integer.  By
Lemma~\ref{lem:v12-anchor-coefficient}, old and new anchors have the
same projection onto all their nonprivate coordinates.  In each
direction use \eqref{eq:v12-quaternion-criterion} to choose the private
witnesses; they need not be the same witnesses in both anchors.
The affine translation of $w$ is invertible.

The exact allowed subgroup is unchanged:
\begin{equation}\label{eq:v12-anchor-subgroup}
 s\in2\,3^{2^{24}-1}\Z\setminus\{0\},\qquad
 \log_2|s|\ge1+(2^{24}-1)\log_2 3>25165823.5.
\end{equation}
Indeed $v_3(f_1)\ge v_3(C)=2^{25}$ is precisely
$2+2v_3(s)\ge2^{25}$.  The subgroup permits arbitrarily divisible
positive $s$, as required by completeness.

\subsection{A shorter scalar norm tree}\label{subsec:v12-scalar}

Keep the exact contact polynomial and finite correction
\[
 r=\beta(\beta-1)(\beta-2401)(32\beta+343)(2\beta-1),\qquad
 c\in\left\{\pm\prod_{p\in I}p:
 I\subseteq\{2,3,5,7,11,13,17\}\right\}.
\]
Here $r\ne0$ for rational $X$.  The only positive contact not excluded
by $\beta\ge3$ is $\beta=2401$, which would imply $X^2=1199$,
impossible over $\Q$.  The old inverse equation for $r$ and its
power chain remain in place.

Write
\[
 f(H)=h_0+h_1\rho,\qquad z(R)/z(H)=z_0+z_1\rho,
 \qquad q=h_0^2-2h_1^2=\operatorname{Norm}_{K/\Q}(f(H)).
\]
Use $N_r(u,v)=u^2+rv^2$ and replace the old six-leaf scalar tree by
\begin{equation}\label{eq:v12-scalar-tree}
 W=N_r\left(N_r(1,q/r),\ N_r(z_0/r^{32},z_1/r^{32})\right).
\end{equation}
The scalar pole input is still exactly $\mathfrak b=crW^2$, with the
original preparation and nonzero condition.  This tree changes that
input's value, so completeness must supply new pole witnesses.

For an explicit quadratic encoding, let $i_{32}=r^{-32}$ be the
retained power-chain output and introduce eight rational coordinates
$q,A,U,V,V_2,B,B_2,W$ with equations
\begin{equation}\label{eq:v12-scalar-eight}
\begin{aligned}
 q&=h_0^2-2h_1^2,& r(A-1)&=q^2,\\
 U&=z_0i_{32},&V&=z_1i_{32},\\
 V_2&=V^2,&B&=U^2+rV_2,\\
 B_2&=B^2,&W&=A^2+rB_2.
\end{aligned}
\end{equation}
The existing $W^2$ coordinate is reattached to this $W$ with its
actual fixed rational normalization.  Thirteen old leaf and norm
coordinates disappear, giving a net saving of five at this stage.
The final elimination of $q$ is given below.

\begin{lemma}\label{lem:v12-scalar-sound}
The new scalar certificate implies
$M_5(\beta)^{32}\mid s$ and
$\operatorname h(\beta)\le\mathcal H|s|^{19/32}$.
\end{lemma}
\begin{proof}
Let $p>76176$ be a prime of the same five-contact radical as in
Lemma~\ref{v5:scalar32-bridge}, and put $\eta=v_p(r)$.
The unchanged contact calculation makes $\eta$ odd; the correction
$c$ is a $p$-unit.  The original pole obstruction, applied to the
retained nonzero input $crW^2$, gives
\[
 \eta+2v_p(W)>0.
\]
At a node $u^2+rv^2$, nonzero summands have opposite valuation parity,
so they cannot cancel.  Zero summands are harmless with valuation
$+\infty$.  If $\eta<0$, the constant-left path in
\eqref{eq:v12-scalar-tree} gives $v_p(A)\le0$ and $v_p(W)\le0$, a
contradiction.  Thus $\eta>0$ and $v_p(W)>-\eta/2$.
Descending the two levels gives
\[
 v_p(A)>-\eta/4,\quad v_p(B)>-3\eta/4,
\]
\[
 v_p(q/r)>-5\eta/8,\quad
 v_p(U)>-3\eta/8,\quad v_p(V)>-7\eta/8.
\]
Therefore
\begin{equation}\label{eq:v12-scalar-valuations}
 v_p(q)>3\eta/8>0,\qquad
 v_p(z_0)>(32-3/8)\eta,\qquad
 v_p(z_1)>(32-7/8)\eta.
\end{equation}
Since $\eta$ is a positive integer, rational integral valuations give
$v_p(z_0),v_p(z_1)\ge32$.

At a split prime $v_p(q)$ is the sum of the two valuations of $f(H)$;
at an inert prime it is twice the single valuation.  Positivity of
$v_p(q)$ therefore supplies at least one place $v$ above $p$ with
$v(f(H))>0$.  We claim only existence of one such place.  At that
place $H$ is formal and $v(z(H))\ge1$.  By the two coefficient bounds,
\[
 v(z(R))\ge v(z(H))+32>0.
\]
This need not make $R$ itself formal: it may lie near $(0,0)$.
We handle both branches explicitly.

Put $m=v(x(R))$.  The retained point domains give $x(R)y(R)\ne0$.
Since $D$ is a unit and
\[
 z(R)^2=\frac{x(R)}{x(R)^2-D},
\]
positive valuation of $z(R)$ excludes $m=0$.  For $m>0$ its valuation
is $m/2$, and for $m<0$ it is $-m/2$.  The exact duplication identity
\begin{equation}\label{eq:v12-duplication}
 x([2]R)=\frac{(x(R)^2+D)^2}{4x(R)(x(R)^2-D)}
\end{equation}
then gives, for either sign of $m$,
\[
 v(x([2]R))=-2v(z(R))<0,\qquad v(z([2]R))=v(z(R)).
\]
Its numerator and denominator are nonzero: with $m\ne0$ their
relevant terms have unequal valuations, and $p$ is odd.
Also $[2]R$ is nonidentity because $s\ne0$ and $P$ has infinite order.

Both $H$ and $[2]R$ are now in the formal subgroup where the logarithm
is a homomorphism and isometry.  Apply it to the exact global identity
\[
 [h]([2]R)=[2s]H.
\]
The points are nontorsion, so their nonzero formal logarithms permit
comparison of valuations.  Since $2$ is a $p$-unit,
\[
 v_p(h)+v(z([2]R))=v_p(s)+v(z(H)),
 \qquad v_p(s/h)\ge32.
\]
The nonzero integer $h$ has nonnegative $p$-valuation, hence
$v_p(s)\ge32$.  This proves the claimed divisibility at every prime
of $M_5(\beta)$.  The same exponent-nineteen height theorem as in
Lemma~\ref{lem:v9-bridge} gives the asserted height bound.
\end{proof}

For completeness use the old scalar preparation with its stronger
six-leaf conditions.  For integral $X$, $r$ is a nonzero integer.
Choose $c$, a sufficiently divisible positive $h_0$, and then a still
deeper positive $s$ satisfying the anchor, exactly as in the proof of
Lemma~\ref{v5:scalar32-bridge}.  At primes dividing $r$, both
coefficients of $f(H)$ are divisible by $r$ locally, so $q/r$ is
integral; the two retained ratio leaves are integral.  At primes
dividing $c$ but not $r$, the same leaves are divisible by the prime.
The constant-left tree is consequently one modulo every prime on
the support of $cr$.  It is nonzero and a unit at those primes.
Thus $crW^2$ is positive, has square reduced denominator, and has the
same local squareclasses as $cr$ at two, three, and seventeen.
The original pole-completeness lemma supplies fresh scalar pole
witnesses.  The new norm and ratio functions are continuous at the
chosen points, so the same finite-neighborhood argument provides a
congruence class of $h_0$ preserving all required conditions.

\subsection{A reciprocal chart for the scalar difference}
\label{subsec:v12-reciprocal}

This replacement is used only for $R=S-P$.  The five comparison
subtractions retain their strengthened guards
\eqref{eq:v12-half-trace}.  Write the incoming scalar point as
$S=(x,hj)$ with $x=h^2$ and $j^2=x^2-D$.  Put
\[
 d=c_*x+D,\quad N=a_*j-b_*h,\quad
 M=-2a_*Dh(x+c_*)-b_*j(c_*x-D).
\]
Both real embeddings give $c_*>0$, $D>0$, and $x\ge0$; hence $d\ne0$.
Introduce $g,k\in K$, retain the existing outputs $f,z\in K$, and
introduce a rational guard $\sigma$.  The final equations are
\begin{equation}\label{eq:v12-reciprocal}
 dg=N,\qquad
 a_*dk+2Dh+a_*b_*j-2Db_*g=0,\qquad
 f=g^2,\qquad f_0\sigma=1,\qquad kz=-g.
\end{equation}
Each field equation means its two rational component equations.
Every component equation is quadratic.  Since
$f_0=g_0^2+2g_1^2$, the guard is precisely $g\ne0$ over rational
coordinates, and imposes no unmentioned sign choice.

We give the identities and all exceptional cases needed to identify
these outputs with the actual group difference.  If
\[
 G=dg-N,\quad K_0=d^2k-M,\quad
 L=a_*dk+2Dh+a_*b_*j-2Db_*g,
\]
then expansion using $a_*^2=c_*$ and $b_*^2=c_*^2-D$ gives
\begin{equation}\label{eq:v12-dcancel}
 dL=a_*K_0-2Db_*G.
\end{equation}
Thus the first two equations of \eqref{eq:v12-reciprocal} are
exactly equivalent to $g=N/d$, $k=M/d^2$, because $a_*d\ne0$.
The latter expressions satisfy
\begin{equation}\label{eq:v12-reciprocal-identities}
 k^2+Dg^4=1,\qquad
 (a_*j-b_*h)(a_*j+b_*h)=(x-c_*)(c_*x+D).
\end{equation}
For example the first is the polynomial identity
$M^2+DN^4=d^4$ modulo $h^2=x$, $j^2=x^2-D$ and the two fixed-constant
identities.  This also specifies an exact denominator-cleared check.

The guard gives $g\ne0$.  If $k=0$, the first identity in
\eqref{eq:v12-reciprocal-identities} gives $d_*g^2=1$ or $-1$.
Taking rational norms gives $7\operatorname{Norm}(g)^2=1$, impossible
for a rational norm.  Thus $k\ne0$, and all displayed outputs are
well defined.  If $x\ne c_*$, use the ordinary subtraction slope
$m=(hj+a_*b_*)/(x-c_*)$.  Substitution and clearing denominators give
\[
 m^2-x-c_*=1/g^2,\qquad
 m(x-1/g^2)-hj=k/g^3.
\]
These are the abscissa and ordinate of $S-P$.  Therefore the retained
outputs are exactly $f(R)=g^2$ and $z(R)=-g/k$.

If $x=c_*$, the incoming chart forces
$h=\pm a_*$ and $j=\pm b_*$.  Equal choices of signs give $S=P$
and $g=0$, excluded by the guard.  The other choices give $S=-P$.
For $h=a_*,j=-b_*$ the explicit values are
\[
 g=\frac{-2a_*b_*}{c_*^2+D},\qquad
 k=\frac{c_*^4-6c_*^2D+D^2}{(c_*^2+D)^2}.
\]
They give the abscissa and ordinate of $-2P$ by the ordinary duplication
formula.  Simultaneously negating $h,j$ negates $g,k$ and leaves
$1/g^2$, $k/g^3$, and $-g/k$ unchanged.  Thus this branch also represents
the true difference, not an arbitrary value of an old zero-over-zero
slope.

In the complete present system that extra branch is impossible.
It would give $R=-2P$ and $v_3(f_1(R))=2$ by
Lemma~\ref{lem:v12-anchor-coefficient}, whereas
\eqref{eq:v12-anchor} requires $v_3(f_1(R))\ge2^{25}$.
It follows that every final solution has $x\ne c_*$.  Restore the old
inverse of $x-c_*$, slope, slope square, and ordinate products by their
ordinary rational formulas.  All retained outputs coincide.  Conversely
every old solution extends to \eqref{eq:v12-reciprocal} by the displayed
formulas, with $\sigma=1/f_0$; its actual nonzero $f(R)=g^2$ ensures
$f_0>0$.  This is an equivalence on all nonprivate coordinates of the
complete system, with the anchor explicitly used to exclude the
additional branch.

For the coordinate count one may first use $d_2=d^2$ and $d_2k=M$.
That intermediate reciprocal chart replaces eight old scalar
difference coordinates by the six components of $g,d_2,k$ and one
rational guard, saving one.  Identity \eqref{eq:v12-dcancel} then removes
the two components of $d_2$.  Forward, divide $dL=0$ by the already
nonzero $d$; backward, reinsert the unique $d_2=d^2$ and use
$a_*K_0=0$.  This second step is a direct rational equivalence requiring
no new exceptional-point argument.  The combined saving is three.

\subsection{A shorter weight circuit with both outputs retained}
\label{subsec:v12-toric}

In each detector let $e$ be its weight $e$, $q=q_1$, and keep both
outputs $q_{12},q_{40}$.  The grouped evaluations already contain
$Q=q^3$ and $T=q_{12}^2$.  Replace the six old intermediate powers of
$e$ by five rational coordinates $\alpha,\gamma,u,v,\tau$ and impose
\begin{equation}\label{eq:v12-toric}
 e\alpha=1,\quad e\gamma=\alpha^2,\quad
 \gamma u=eq,\quad \gamma v=u^2,\quad
 Qq_{12}=v^2,\quad q\tau=eT,\quad vq_{40}=T\tau.
\end{equation}
These are seven quadratic equations for seven coordinates when the
two retained outputs are counted.  The old circuit had eight.

The first equation gives $e\ne0$.  The unchanged definition of $q$
is $q\delta^3=e$, with the powers of $\delta$ represented by their
retained coordinates, so $q\ne0$ as well.  The equations now determine,
in order,
\[
 \alpha=e^{-1},\quad\gamma=e^{-3},\quad
 u=qe^4,\quad v=q^2e^{11},\quad q_{12}=qe^{22}.
\]
The division by $Q=q^3$ is legitimate.  The retained equation for
$T$ then gives $T=q^2e^{44}$, and the last two equations give
\[
 \tau=qe^{45},\qquad q_{40}=qe^{78}.
\]
In particular $v\ne0$ whenever it is used as a divisor.  This is a
successive computation, not a circular specification involving $T$.

The forward map from the old circuit assigns these five Laurent
monomials and retains $q_{12},q_{40}$, $Q,T$, and every other coordinate.
For the reverse map restore the old six powers as
$e^{-1},e^{-3},e^{-7},e^{-14},e^{-29},e^{-55}$.  The old output equations
then give the same $qe^{22}$ and $qe^{78}$.  The replaced internal
powers have no outside consumers.  Thus these maps preserve every
outside coordinate, not merely the input.  Applying this equivalence
to all three detectors saves three coordinates.

\subsection{Removing the scalar norm coordinate}\label{subsec:v12-norm-elim}

One further elimination uses coordinates already present in the
order-forty test $f(H)$.  With $f(H)=h_0+h_1\rho$, its square is stored
through
\[
 t=h_0^2+2h_1^2,\qquad u=h_0h_1.
\]
Here $u$ is half the $\rho$-coefficient, a normalization that matters.
The scalar coordinate $q=h_0^2-2h_1^2$ from
\eqref{eq:v12-scalar-eight} has only one other consumer, the equation
$r(A-1)=q^2$.  Use the identity
\begin{equation}\label{eq:v12-norm-square}
 (h_0^2-2h_1^2)^2=(h_0^2+2h_1^2)^2-8(h_0h_1)^2=t^2-8u^2.
\end{equation}
Delete $q$ and its defining row and replace the consumer by
\[
 r(A-1)=t^2-8u^2.
\]
This is still quadratic.  Forward, forget $q$.  Backward, reinsert
$q=h_0^2-2h_1^2$.  Identity \eqref{eq:v12-norm-square} proves both
rows in the old system and makes these maps mutually inverse.
No division, nonzero hypothesis, or extra arithmetic input is used.
The proof of Lemma~\ref{lem:v12-scalar-sound} applies by this unique
reinsertion; it does not require retaining $q$ as a quantified variable.

\subsection{Soundness of the combined construction}
\label{subsec:v12-soundness}

Let a rational solution of the final system be given.  The inherited
odd charts identify the four root points with odd integer labels.
The reciprocal chart supplies the actual $R=S-P$, is nonidentity,
and, together with the anchor, excludes $S=-P$.  The explicit maps
above therefore restore all unchanged point domains and quotient
guards needed by the earlier arguments.  The removed scalar norm
coordinate also has its unique rational value.

Put $\kappa=|s|=|a-1|$.  By
\eqref{eq:v12-anchor-subgroup} and Lemma~\ref{lem:v12-scalar-sound},
\[
 \log_2\kappa>25165823.5,\qquad
 \operatorname h(\beta)\le\mathcal H\kappa^{19/32}.
\]
These are the exact bounds consumed in the Version 9 height argument.
For completeness we repeat the numerical deductions.  Since
\[
 10(7+2093966)<25165823.5,\qquad80<2^7,
\]
we have
\[
 \kappa>(80\mathcal H)^{10},\qquad
 |a|\ge\kappa-1\ge\kappa/2>(40\mathcal H)^{10},
\]
and
\begin{equation}\label{eq:v12-margins}
 \operatorname h(\beta)<\kappa/80\le|a|/40,\qquad
 \log(40\mathcal H)<\kappa/80\le|a|/40.
\end{equation}
The first inequality uses
$\kappa^{13/32}>80\mathcal H$; the logarithmic inequality is weaker.
In particular all labels are in the large range needed below.

Set
\[
 F=t^2-a(a-2)tb+(a-2)^2b^2-1,\qquad
 B=\max(|n|,|a|,|t|,|b|)>4.
\]
The triangle inequality gives $|F|<5B^4$.  If $F\ne0$,
Lemma~\ref{lem:v12-mixed-local} and the unchanged detector implication
make the three contact radicals divide $F$ to order forty.
Equation~\eqref{eq:v12-height-input} then gives
\[
 B^{1/10}<7\mathcal H\,5^{19/40}+6<24\mathcal H,
\]
contradicting $B\ge|a|>(40\mathcal H)^{10}$.  Hence $F=0$.
The point $Q_4$ is nonidentity, so its integer label $b\ne0$.
Lemma~\ref{lem:o8-lucas} therefore gives $\log B>|a|-3$.

If $n\ne\beta$, let $N$ be the nonzero reduced numerator of
$n-\beta$.  Then $|N|\le2e^{\operatorname h(\beta)}B$, and the same
three radicals divide $N$ to order twelve.  Using
$B^2-6>(5/8)B^2$ for $B>4$, equation~\eqref{eq:v12-height-input}
gives
\[
 B^{5/12}<40\mathcal H e^{(19/12)\operatorname h(\beta)}.
\]
Consequently
\[
 \log B<\frac{12}{5}\log(40\mathcal H)
            +\frac{19}{5}\operatorname h(\beta)
       <\frac{31}{200}|a|<|a|-3,
\]
a contradiction.  Thus $n=\beta\in\Z$.  If $X=u/v$ in lowest terms
with $v>0$, integrality of $2X^2+3$ forces $v^2\mid2$, so $v=1$.
Every rational solution therefore projects to an integer input.

\subsection{Completeness with all replacements at once}
\label{subsec:v12-completeness}

Fix an integer $X$ and put $n=\beta=2X^2+3$, a positive odd integer.
First perform the scalar preparation with the original six-leaf
bounds, which, as shown above, suffice for the new scalar tree.
Choose its finite correction and a sufficiently divisible positive
shared index $h_0$, with $64\mid h_0$.  Choose a positive scalar index
$s$ still more divisible, large enough for the physical anchor, and
with $6\mid s$.  Put $a=s+1$.  The actual prepared scalar point is
$R=[s]P=S-P$.  The new anchor has exactly the old allowed subgroup,
and the new scalar tree has the required unit root.  Fresh scalar
pole witnesses therefore exist.

To choose the two other root labels, define
\[
 U_0=0,\quad U_1=1,\quad U_{j+1}=aU_j-U_{j-1},\qquad
 k=2(a-2),\quad b=U_k/(a-2),\quad t=U_{k+1}.
\]
Induction gives $U_j\equiv j\pmod{a-2}$, so $b$ is an integer.
Modulo two the sequence has period $0,1,1$.  Since $a\equiv1\pmod3$,
$k\equiv1\pmod3$, and both $U_k,U_{k+1}$ are odd.  The positive odd
denominator $a-2$ makes $b$ odd.  The Lucas identity gives
$F(a,t,b)=0$.  Thus all four root indices are positive and odd, and
Proposition~\ref{v9:prop:odd-chart} supplies their chart witnesses.

Fix these root points and the three positive rational square scales
of Corollary~\ref{v9:cor:scales}.  The new scalar functions and the
mixed alternate numerator are rational functions defined at the
prepared points.  Their required finite local conditions persist in
a sufficiently small neighborhood.  As in the original argument,
choose a common continuity modulus $L$, with $64\mid L$, so that
$h\equiv h_0\pmod L$ preserves them.  All these choices precede the
contact-prime selection.  Increase the finite prime cutoff to include
the fixed root denominators needed for the five half-trace guards.

Choose the contact primes using the same simultaneous-prime theorem
and Lemma~\ref{v9:lem:separation}.  Put
\[
 M_\pm=\operatorname{lcm}_p(d_\pm(p)/4),\qquad\Pi=\prod_p p.
\]
The retained direct-target construction supplies
\[
 \gcd(M_-,M_+)=1,\quad\gcd(M_-,L)=16,\quad\gcd(M_+,L)=1,
 \quad\gcd(\Pi,M_+)=\gcd(\Pi,L)=1.
\]
Thus the same three congruences are compatible:
\[
 h\equiv0\pmod{\Pi^{19}M_-},\qquad
 h\equiv1\pmod{M_+},\qquad h\equiv h_0\pmod L.
\]
Use the inherited real-progression argument to choose a sufficiently
large positive $h$ in this progression, avoiding the finitely many
forbidden point and quotient loci.  Set
$U_i=[hn_i]P$ and $V_i=[(h-1)n_i]P$.

At every primary place the points $H,U_i$ have depth at least twenty:
the doubling-image order first makes them formal, and multiplication
by $p^{19}$ raises depth by nineteen under the same formal logarithm
isometry.  At opposite places the points $V_i$ are formal.  Hence
$L_i$ agrees with its integer label to order at least eighty; the
alternate ratio agrees to order twelve for $i=1$ and forty for the
other three labels.  All tests \eqref{eq:v12-tests} hold.  The same
primary/opposite reductions prove that all five half-traces are
nonzero, by the argument following \eqref{eq:v12-half-trace}.
The chosen scalar label has $S\ne\pm P$, so its reciprocal chart
has the ordinary rational lift described above.

The pole-realization estimates are valid for these fixed orders
$1,12,40$.  In the source notation their large-real-parameter decay
is still $\lambda^{-3}$; the small-prime lower bound is
$3\nu+(10m-8)\eta-2c$, which tends to infinity under the same
preparatory divisibility.  No new prime-pattern assertion is used.
Obtain the detector pole witnesses and apply the rational
normalization \eqref{eq:v12-pole-normalized}.  Apply the weight-circuit
maps \eqref{eq:v12-toric}, the reciprocal denominator cancellation,
and the norm-square elimination.  These exact maps preserve all
outside coordinates.  They therefore give one simultaneous rational
solution of the final system for the chosen integer $X$.

This proves completeness in the required order of choices.  In
particular, it does not combine separately chosen primes or
incompatible continuity moduli from different optimizations.

\subsection{Exact polynomial, counts, and coefficient lift}
\label{subsec:v12-finite}

Each equation above is represented by its rational component rows;
intermediate products in the retained straight-line program are
represented by their specified coordinates.  Normalize each nonzero
rational residual by clearing denominators and dividing its integer
content, and let $P_{498}$ be the sum of the squares of those 467
normalized residuals.  This is an integer polynomial.  A sum of
rational squares vanishes exactly when every residual vanishes.
The soundness and completeness just proved therefore establish
Theorem~\ref{thm:v12-quartic}.

The following count records the actual transported, fully reduced
systems; it does not count unreduced formulas in isolation.
\begin{center}\small
\begin{tabular}{lrr}
\toprule
System or replacement & Total variables & Quadratic residuals\\
\midrule
The retained Version 9 baseline &523&492\\
Mixed input precision and its checked compilation &519&488\\
Three normalized detector poles &516&485\\
Five half-trace guards &511&480\\
Single-coefficient translated anchor &510&479\\
Shorter scalar norm tree &505&474\\
Scalar reciprocal chart with $d^2$ retained &504&473\\
Three weight circuits retaining $q_{12},q_{40}$ &501&470\\
Cancellation of the scalar $d^2$ coordinates &499&468\\
Norm-square coordinate elimination &498&467\\
\bottomrule
\end{tabular}
\end{center}

The mixed-depth compilation uses only exact quadratic identities,
constant-pivot substitutions, and invertible shears in shared
coefficient coordinates.  A constant-pivot substitution inserts its
unique polynomial value in the reverse direction; a shear inserts
its inverse linear transformation.  Both maps fix the input.
The final norm-square elimination is the explicit identity
\eqref{eq:v12-norm-square}.  Thus no compiler reduction introduces a
variable denominator or alters an arithmetic guard.  The exact
certificates also check that private coordinates deleted in the
other replacements have no undeclared consumers.  In particular,
the half-trace witnesses occur only in their guard equations; the
removed weight powers occur only in their old weight circuit; the
removed scalar $d^2$ coordinates occur only in their definitions and
$k$ equations; and the removed norm occurs only in its definition and
constant-node equation.

The frozen residual file is \texttt{candidate498-ir.json}, with
SHA-256
\begin{center}\small
\nolinkurl{7682a1229149c75e27322e86dec424449b55321b06c48e468f90fa78bcebb61a}.
\end{center}
The complete sparse coefficient table \texttt{polynomial498.json}
has SHA-256
\begin{center}\small
\nolinkurl{817fec936dee3b358fc8b5b0b9a5570694adf009b134432467ef0be914aeedc4}.
\end{center}
Its exact expansion has 24,669 monomials and degree four; all 498
variables occur.  A coefficient of absolute value one proves content
one.  The only quadratic-in-$C$ coefficient is $16C^2$, at the square
of the translated anchor witness $w$, with no lower-$C$ term at that
monomial.  All other coefficients are independent of $C$ or linear
in $C$ and are smaller in absolute value.  The coefficient comparison
is carried out symbolically; no expansion of $3^{2^{25}}$ is needed.
It gives the exact maximum $16C^2=16\cdot3^{2^{26}}$.

There is also an ordinary-coefficient lift with sixteen additional
coordinates:
\[
 c_0=3^{1024},\qquad c_{j+1}=c_j^2\quad(0\le j<15),\qquad C=c_{15}.
\]
The recurrence has exactly one rational solution and ends at
$3^{1024\cdot2^{15}}=3^{2^{25}}$.  Replace $-4Cw$ in
\eqref{eq:v12-anchor} by $-4c_{15}w$.  All residuals remain quadratic.
This gives a specified quartic in 514 variables with the same input
projection and 513 existential witnesses.  It has 483 quadratic
residuals and 24,716 monomials, with content one and maximum coefficient
$3^{2048}+1005072314676622<2^{3247}$.  Its coefficient table,
\texttt{polynomial514.json}, has SHA-256
\begin{center}\small
\nolinkurl{d247a27709afda5f2be626390cffa1a175c3dfcbee1a5c832a1c457b32629bd5}.
\end{center}
The lift is unique, so it requires no additional arithmetic argument.

The exact checks establish the residual identities, the witness
counts, and every coefficient of the displayed sums of squares.
They are separate from the arithmetic arguments above.  The
inherited source qualifications of Section~\ref{subsec:v9-inputs}
remain in force; a finite polynomial check is not a new proof of the
restricted converse, the height theorem, or the simultaneous-prime
input.
